% 【本文件写什么】api-v0 契约层，叶子，只写语义不写实现
% - crate 路径：os/components/wateros-mm/mm-api/api-v0/src/
% - 列出并说明：pub trait、核心类型、错误枚举、常量
% - 每个公开 API 的前置条件、后置条件、线程/中断上下文限制
% - 与 Linux/用户态约定的对齐点与 deliberate 差异
% - v0 版本当前行为与后续可替换 extension 点
% 【事实来源】
% - 上述 crate 下全部 .rs 源文件
% - docs/exports/public-api/ 中对应符号表，以源码为准
% 【不写什么】
% - impl 内部数据结构、汇编、具体算法步骤

\subsubsection{地址、权限与错误}

\code{addr}模块提供\code{VirtAddr}、\code{PhysAddr}、\code{VirtPageNum}、\code{PhysPageNum}与常量\code{PAGE\_SIZE}。虚拟/物理地址为字节粒度，页号由\code{floor\_page}/\code{ceil\_page}按 4\,KiB 对齐分解。\code{perm::PagePerm}表达 R/W/X/U 组合；\code{flags::MapFlags}承载\code{MAP\_PRIVATE}、\code{MAP\_SHARED}、\code{MAP\_FIXED}等 mmap 语义子集。\code{error::MmError}统一\code{OutOfMemory}、\code{InvalidAddress}、\code{AlreadyMapped}、\code{NotMapped}、\code{AccessViolation}、\code{Unsupported}及来自帧分配器的\code{FrameAlloc}嵌套错误。

\subsubsection{地址空间与映射契约}

\code{AddressSpaceOps}是页表操作的语义核心：\code{map\_page\_to\_ppn}、\code{unmap\_page\_to\_ppn}、\code{protect\_page}、\code{translate\_addr}、\code{satp\_value}。默认方法\code{map\_page\_with\_alloc}委托\code{PhysicalFrameAllocator}分配帧后映射。\code{fork}默认返回\code{Unsupported}，真实 fork 由\code{mm-impl}覆盖。

\code{MmapOps}在\code{AddressSpaceOps}之上扩展\code{mmap}、\code{munmap}、\code{mprotect}、\code{mremap}，Linux 语义子集、\code{handle\_page\_fault}，惰性文件/匿名缺页。\code{MmapRequest}携带\code{addr\_hint}、\code{len}、\code{prot}、\code{flags}、\code{MmapKind::Anonymous | File}。\code{DemandPageLoader} trait 供文件懒映射与 fork 后 loader 复制。

\begin{rustcode}
pub trait MmapOps: AddressSpaceOps {
    fn mmap<A: PhysicalFrameAllocator<FrameId = PhysPageNum>>(
        &mut self, allocator: &mut A, req: MmapRequest,
        file_backing: Option<&[u8]>,
    ) -> MmResult<VirtAddr>;
    fn handle_page_fault<A>(
        &mut self, allocator: &mut A,
        fault_addr: VirtAddr, access: PageFaultAccess,
    ) -> MmResult<bool>
    where A: PhysicalFrameAllocator<FrameId = PhysPageNum>;
}
\end{rustcode}

\subsubsection{堆、用户访问与 bring-up}

\code{brk::HeapBrk}描述堆区间\code{BrkRegion}与\code{brk(new\_end)}语义：跨越新页时须分配并映射整页，默认权限 R/W/U。\code{user\_access::UserMemoryOps}定义 syscall 路径对用户缓冲区的读写字节序列与\code{probe}诊断接口，须在当前任务已安装的用户地址空间上下文中调用。

\code{kernel\_bringup}定义\code{LoadedElf}，含入口 PC、映像区间、栈界、\code{user\_aspace\_ptr}、\code{satp} token，以及 \code{LoadElfError}、\code{LoadProgramError}，含 shebang 解析。\code{elf\_user\_stack}按 RISC-V 用户栈布局写入 argc/argv/envp/auxv。\code{kernel\_satp}缓存内核地址空间 token。\code{user\_aspace\_lifecycle}注册任务 exit 时释放用户地址空间的钩子。\code{mempolicy}为单节点 NUMA 桩，满足\code{get\_mempolicy}返回值约定。

\code{PhysicalFrameAllocator} trait 定义于本 crate 的\code{frame\_allocator}模块，帧标识类型与\code{PhysPageNum}对齐；具体分配器在\code{frame-alloctor}子系统中实现。
